\documentclass{article}
\usepackage{graphicx}

\title{Towards Automated Goal Driven Dataset Discovery}

\begin{document}

\maketitle

\section{Introduction}
Difference between harini and vidits methods for mining datasets and the principal architectures that they adopted because of audience vs phenomenon. This provides an idea that we can pivot to one or more of lasswel's pillars. 

While earlier researchers used to choose simpler or broad or specific methodology that caters to the head of the distribution because this gives them the best advantage to human resources ratio, however this cost of data research is amortized among agents.

But this limits the kind of phenomenon, audiences, and contexts thus curtailing the research to qualitative aspects only. For e.g How to persuade italian gamblers. Subreddit specific comments and posts, data mining, dataset construction which couldnt have been done by humans earlier, can now be automated, increasing the total datasets and studied phenomenon and different aspects of it.

Things that required human annotation like memorability

We can choose the subreddit to model any of the dimensions of lasswel.
Subreddit and the specific behaviors observed over there -> Harini made audience as the pivot while selecting the phenomenon (top 20 vs remaining authors) whereas Vidit made phenomena as the pivot without imposing audience. Harini could make CMV as the pivot with self attested deltas. 

Subpop -> Audience (Demographic) x Surveys -> economic, social

% BehaVoir — Scratch Pad
% Dump ideas, half-thoughts, references, tidbits here.
% When an idea matures, move it into the main paper and mark it here.
%
% Format:
%   \idea{Short title}{Details or stream of consciousness}
%   \ref{What to look up}{Why it matters}
%   \snippet{Raw text or formula to possibly include}

\newcommand{\idea}[2]{\paragraph{Idea: #1} #2}
\newcommand{\reference}[2]{\paragraph{Ref: #1} #2}
\newcommand{\snippet}[1]{\begin{quote}\textit{#1}\end{quote}}

%  IDEAS
\section{Ideas}

\idea{Reddit as primary data source}{Why Reddit? (1)~Data availability --- public API, Pushshift archives make large-scale collection feasible. (2)~Rich structure --- subreddits act as natural ``environments,'' plus threaded conversations, voting signals, user histories, and timestamps. (3)~Massive scale --- billions of posts/comments spanning years of organic human behavior. (4)~A moving snapshot of mankind --- ongoing, real-time record of behavior across thousands of topics and communities. Key framing: Reddit's subreddit structure naturally partitions human behavior into discoverable ``environments'' --- exactly what BehaVoir needs for dynamic dataset discovery. Note: demographic biases (age, gender, geography skew) are acknowledged but out of scope.}

\section{References}
% References
\begin{itemize}
    \item Our work currently like AlphaGo: Experigen \cite{gupta_accelerating_2026}
    \item Our next steps are to extend this for transfer learning and multi task hypothesis generation similar to \cite{silver_mastering_2017}
    \item Dreamer/MuZero \cite{schrittwieser_mastering_2020}: What we want to achieve for human behavior\cite{hafner_dream_2020}
    \item Lasswell's communication model (1948): \cite{lasswell1948structure} --- theoretical backbone for our Reddit mapping
    \item Reddit social roles / community structure: \cite{buntain2014identifying}
\end{itemize}



% --- 2026-05-12: Core Theoretical Framing ---

\idea{Communication as the sole learning signal}{%
We disassociate from cognition entirely. Everything can be learned from human
\emph{actions} --- i.e., observations. Communication IS the action.
Data is available and as complete as we can wish for at the level of
individuals, communities, and social structures --- to the point individuals
express themselves through communication. Humans themselves only need to
``see'' or observe each other's actions to understand each other, hence
observation should be sufficient. This \textbf{jagged intelligence} is what
we explore and optimize.}

\idea{Lasswell's Model as theoretical backbone}{%
Communication, from Lasswell's model:
\textbf{Who} $\to$ \textbf{Says What} $\to$ \textbf{In What Channel}
$\to$ \textbf{To Whom} $\to$ \textbf{When} $\to$ \textbf{With What Effect}.
\\[4pt]
Reddit mapping:
\begin{itemize}
    \item \textbf{Who / To Whom}: Imperfect --- individuals may hide or fake identity, but that is the true nature of humans. Assumed \emph{nearly faithful}.
    \item \textbf{What + Context (Channel)}: This is where ALL Reddit research focuses today --- researchers manually identify structures (subreddits, rules) to give structure to observations (CMV, WritingPrompts, AITA). Even identity is expressed through communication (MBTI signals, etc.). \textbf{We want to automate this and use compute to reach superhuman performance.}
    \item \textbf{Effect}: We don't have individual atomic upvotes, but aggregate upvotes exist. \textbf{Comments are the atomic unit of conversation.}
    \item \textbf{When}: Reddit is a moving time series $\Rightarrow$ enables \textbf{pre-registered experiments}, a gold standard in observational sciences. This completes the feedback loop.
\end{itemize}}

\idea{LLMs as the engine}{%
LLMs serve as:
(1)~A general reservoir of \emph{noisy world knowledge},
(2)~A compute-efficient automated reasoner,
(3)~Capable at code $\Rightarrow$ scalably voyaging through the labyrinth of human behavior.
\\[4pt]
Mapping LLMs to Lasswell:
\begin{itemize}
    \item They model \textbf{``what''} which in turn captures ``who, what, whom'' decently enough
    \item They \emph{don't} know how to interleave it in scientific contexts $\Rightarrow$ weak on \textbf{``channel''}
    \item Have a diffused sense of \textbf{``when''} (temporal reasoning is weak)
    \item Starting to show decent prior in \textbf{``effect''}: Hewitt et al.\ (2024) show GPT-4 predicts treatment effects at $r{=}0.85$ vs.\ actual outcomes across 70 pre-registered experiments. See also \cite{park2023generative}, \cite{horton2023homo}, \cite{manning2024automated}, \cite{aher2023using}, \cite{argyle2023out}, \cite{santurkar2023whose}
\end{itemize}
The gap: LLMs have the pieces but lack the scientific scaffolding to compose them. BehaVoir provides that scaffolding.}

% --- 2026-05-12: MuZero Analogy Refinements ---

\idea{MuZero's real success factors and what they map to}{%
MuZero works \emph{not} because its learned model is accurate --- it's objectively
inaccurate and degrades over longer rollouts. It works because:
\begin{enumerate}
    \item \textbf{Value-equivalent modeling}: it only predicts what matters for
          \emph{decisions} (reward, value, policy), not pixel-perfect reconstruction.
          $\Rightarrow$ BehaVoir only needs to discover what's \emph{scientifically
          relevant}, not perfectly structure all of Reddit.
    \item \textbf{The policy prior regularizes search}: the prior steers MCTS toward
          regions where the learned model is \emph{more accurate}. Without it, planning
          in complex environments (Atari) fails completely.
          $\Rightarrow$ \textbf{LLM priors are not just convenient --- they may be
          necessary.} The LLM's existing world knowledge regularizes BehaVoir's
          environment search toward scientifically meaningful regions. Without this
          prior, raw Reddit is too noisy (the ``Atari without policy prior'' problem).
    \item \textbf{Planning self-corrects model errors}: short lookahead in
          approximately-right latent space suffices; MCTS compensates for dynamics
          errors.
          $\Rightarrow$ Reddit's time-series nature enables \textbf{pre-registered
          experiments} = our MCTS. Discovered environments don't need to be perfect;
          the empirical testing loop corrects errors in the learned model.
    \item \textbf{End-to-end training}: representation, dynamics, and prediction are
          jointly optimized. The latent space is whatever makes prediction accurate,
          not ``the true state of the world.''
          $\Rightarrow$ BehaVoir's pipeline should be end-to-end: environment
          discovery, hypothesis generation, and validation jointly optimized.
\end{enumerate}}

\idea{MuZero caveats that constrain the analogy}{%
\begin{itemize}
    \item \textbf{Model is local, not universal}: MuZero's dynamics are only accurate
          near its own data-collection policy. Implication: discovered environments
          may be \emph{local} to the current scientific question, not a universal map.
          Could be a feature (targeted discovery) --- but be explicit.
    \item \textbf{Sample efficiency was the Achilles heel}: vanilla MuZero needed
          $\sim$200M frames. EfficientZero fixed this with self-supervised consistency.
          Our equivalent: how much raw data before the environment model is useful?
          Self-supervised consistency (does a discovered environment stay coherent
          over time?) may be essential.
    \item \textbf{Planning depth limited by model accuracy}: MuZero plateaus at
          $\sim$100 simulations in messy environments. For BehaVoir: near-term
          hypotheses are tractable; long causal chains are hard. Lasswell's ``with
          what effect'' is the hardest dimension --- it requires the longest horizon.
    \item \textbf{Opacity}: MuZero's latent states are uninterpretable. If BehaVoir
          uses learned representations, we must address interpretability or
          accept a similar tradeoff.
\end{itemize}}

\idea{The deepest lesson from MuZero}{%
The bar is not ``discover the perfect dataset.'' It's ``discover datasets that are
\emph{good enough} to generate hypotheses that self-correct through empirical testing.''
The pre-registered experiment capability (Reddit as time series) is our planning
mechanism --- it's how we verify and correct the learned model. Accuracy of discovery
matters less than the \emph{closed loop}: discover $\to$ hypothesize $\to$ test $\to$
refine. MuZero's model is wrong everywhere, but \emph{less wrong where it matters}.}

% --- Post-MuZero research trajectory (2021--2024) ---
\idea{Post-MuZero lineage}{%
Every major follow-up attacked a specific limitation, not the core architecture
(learned model + tree search + policy prior held throughout):
\begin{itemize}
    \item EfficientZero (2021): sample efficiency via self-supervised consistency
    \item Sampled MuZero (2021): continuous/high-dimensional action spaces
    \item Stochastic MuZero (2021): stochastic environments
    \item MuZero Unplugged (2021): offline learning from logged data
    \item Gumbel MuZero (2022): planning with very few simulations ($\sim$2!)
    \item UniZero (2024): long-term dependencies via transformer world models
    \item EfficientZero V2 (2024): unified discrete + continuous, dense + sparse
\end{itemize}
Pattern: the skeleton survived; each piece got more robust. For BehaVoir,
this suggests getting the \emph{skeleton} right matters more than perfecting
any single component.}

% --- 2026-05-12: Field Naming Decision ---

\idea{Field name: Computational Social Science}{%
Recommendation from literature survey:
\begin{enumerate}
    \item \textbf{Computational Social Science (CSS)} --- dominant term at target venues
          (ICWSM, WWW, CHI, ACL). EXPERIGEN already uses ``social science research.''
          Chris Bail's Oxford series, Diyi Yang's Stanford keynotes all rally under this.
          \textbf{Use as primary field label.}
    \item \textbf{Computational Behavioral Science} --- useful to narrow scope to
          \emph{behavior} vs.\ macro social phenomena (elections, policy, economics).
          Not yet an established venue/track name. Good as a positioning phrase within
          the abstract.
    \item \textbf{AI for Social Science} --- gaining traction at AAAI, NeurIPS, ICML
          workshops. Foregrounds the AI contribution. Use as secondary framing for
          ML audience.
\end{enumerate}
Strategy: lead with \textbf{``computational social science''} as the field, narrow in
intro by saying we focus on understanding human behavior from observational
communication data. Use ``AI for social science'' when addressing ML audience.}

% --- 2026-05-12: Curated Observational Datasets Catalog ---

\idea{Reddit-based research datasets (evidence for the ``clean data bottleneck'')}{%
Social science progress is bottlenecked by clean observational datasets.
Researchers have spent years curating a small number of Reddit sources because
the scientific method needs clean formulation, high reward signals, and clear
context. Using all of Reddit is too noisy. Below: the datasets that survived
this curation filter.

\textbf{Reddit-Based Datasets:}
\begin{enumerate}
    \item \textbf{r/ChangeMyView --- Persuasion.}
          Delta mechanism gives binary ``mind changed'' label.
          Structured debate: OP states view, challengers argue, delta = ground truth.
          Tan et al.\ (2016), ``Winning Arguments'' (Cornell). ConvoKit corpus:
          34,911 speakers, 293K utterances. The gold standard for computational persuasion.

    \item \textbf{r/AmItheAsshole --- Moral Judgment.}
          Community votes YTA/NTA/ESH/NAH give distributional ethical labels.
          Lourie et al.\ (2021), \textsc{Scruples} (AllenAI/AAAI): 32K anecdotes,
          625K judgments. Forbes et al.\ (2020), Social Chemistry 101. Clean because
          the voting format forces a verdict on a real-world dilemma.

    \item \textbf{r/WritingPrompts --- Creative Generation.}
          Short prompt $\to$ long story pairs. Fan et al.\ (2018), ``Hierarchical
          Neural Story Generation'' (Facebook/ACL): 300K stories, 97K prompts.
          Clean because prompt--response structure is a natural seq2seq task.

    \item \textbf{r/SuicideWatch --- Suicidal Ideation / Mental Health.}
          Self-disclosure in a support community. Shing et al.\ (2018), UMD Reddit
          Suicidality Dataset; CLPsych shared tasks (2019, 2024). Expert-annotated
          risk on a 4-point scale. Clean because the subreddit self-selects a
          population expressing genuine distress.

    \item \textbf{Mental health subreddits (r/depression, r/anxiety, r/bipolarreddit, etc.) --- Psychiatric NLP.}
          15+ support-group subreddits. Low et al.\ (2020), COVID mental health
          study (JMIR). Benchmark suite: Zenodo dataset (28 subreddits).
          Clean because each subreddit maps to a diagnostic category.

    \item \textbf{r/explainlikeimfive + r/askscience + r/AskHistorians --- Long-Form QA.}
          ELI5 dataset (Fan et al., 2019, Facebook): question--answer pairs
          requiring paragraph-length explanation. Clean because community norms
          enforce sourced, accessible answers; upvotes reward quality.

    \item \textbf{Self-Annotated Reddit Corpus (SARC) --- Sarcasm Detection.}
          Users' own \texttt{/s} tag = self-labeled sarcasm. Khodak et al.\ (2018,
          LREC/Princeton): 1.3M sarcastic statements with conversational context.
          Clean because annotation is author-provided, not crowdsourced.

    \item \textbf{r/TheOnion vs.\ r/nottheonion --- Satire/Sincerity Detection.}
          Parallel subreddits: one satirical, one real-but-absurd news. Natural
          binary classification task. Used in multiple NLP classroom + research
          projects ($\sim$80--87\% accuracy). Clean because the subreddit boundary
          is the label.

    \item \textbf{Moral Foundations Reddit Corpus (MFRC) --- Moral Sentiment.}
          16,123 comments from 12 subreddits (politics, confessions, AITA, etc.),
          hand-annotated for 8 moral-foundation categories. Trager et al.\ (2022).
          Clean because trained annotators applied a validated psychology codebook.

    \item \textbf{r/confession / r/offmychest --- Regret and Confessional Behavior.}
          Prabhala et al.\ (2023), ReDDIT dataset: regret detection (action vs.\
          inaction vs.\ no regret). Also used in Judge-BERT moral-judgment study.
          Clean because users self-select to confess norm violations.

    \item \textbf{r/The\_Donald (quarantine studies) --- Moderation \& Radicalization.}
          Chandrasekharan et al.\ (2022), ``Quarantined!'' (ACM TOCHI): 85M posts,
          causal inference on quarantine effects. Clean because the policy intervention
          (quarantine date) provides a natural experiment with a sharp treatment boundary.

    \item \textbf{r/Liberal vs.\ r/Conservative --- Political Ideology / Media Bias.}
          226K shared articles annotated for ideological bias. Baly et al.\ (2024,
          ScienceDirect). Clean because subreddit membership proxies ideology;
          shared URLs link to external ground truth.
\end{enumerate}

\textbf{Non-Reddit Observational Datasets (same pattern):}
\begin{enumerate}
    \item \textbf{Wikipedia Talk Pages --- Toxicity.}
          Jigsaw/WikiDetox: 160K comments, 6 toxicity labels; 100K diffs with 1M
          crowdsourced annotations. Wulczyn et al.\ (2017), ``Ex Machina'' (WWW).
          Perspective API. Clean because talk pages have norms + extensive moderation logs.

    \item \textbf{Twitter/X --- Sentiment \& Hate Speech.}
          SemEval shared tasks (2013--2019): sentiment (Rosenthal et al., 2017),
          hate speech (Basile et al., 2019, HatEval), offensive language (Davidson
          et al., 2017). TweetEval benchmark. Clean because shared-task curation
          enforces annotation standards and train/test discipline.

    \item \textbf{Stack Overflow --- Technical QA \& Community Dynamics.}
          SNAP temporal network; annual Developer Survey (49K+ respondents);
          HuggingFace stack-exchange-preferences (RLHF training). 83M+ Q\&A pairs
          over 17 years. Clean because accepted-answer + voting gives explicit
          quality signal.
\end{enumerate}

\textbf{Key takeaway for BehaVoir:} Every dataset above exists because humans
manually identified a subreddit/platform whose \emph{structure} provides free
labels (deltas, votes, \texttt{/s} tags, subreddit boundaries, policy dates).
This is what BehaVoir aims to automate: discovering these clean observational
environments at scale rather than one-at-a-time over decades of manual curation.}

% --- 2026-05-12: LLM Behavior Simulation Literature ---

\idea{LLMs have decent priors on behavioral ``effect'' --- the evidence}{%
Key papers establishing that LLMs encode usable behavioral priors:
\begin{enumerate}
    \item \textbf{Generative Agents} \cite{park2023generative} (UIST 2023): 25 LLM agents in a sandbox autonomously form relationships, spread information, coordinate. Emergent social dynamics without explicit programming.
    \item \textbf{Homo Silicus} \cite{horton2023homo} (NBER 2023): LLMs replicate classic economic experiments qualitatively. Coins the term. Argues LLMs are implicit computational models of humans.
    \item \textbf{Automated Social Science} \cite{manning2024automated} (NBER 2024): LLMs both generate hypotheses AND serve as subjects. Predict \emph{sign} of causal effects well, but not magnitudes. Key: simulation reveals knowledge LLMs cannot directly articulate.
    \item \textbf{Turing Experiments} \cite{aher2023using} (ICML 2023): Replicate Ultimatum Game, Milgram, Garden Path. ``Hyper-accuracy distortion'' --- decent but imperfect.
    \item \textbf{Silicon Samples} \cite{argyle2023out} (Political Analysis 2023): LLMs conditioned on demographics approximate ANES survey distributions.
    \item \textbf{Whose Opinions?} \cite{santurkar2023whose} (ICML 2023): Systematic bias toward college-educated viewpoints. Priors are decent but skewed.
    \item \textbf{Hewitt et al.\ (2024, preprint)}: GPT-4 predicts treatment effects at $r{=}0.85$ across 70 pre-registered experiments (476 effects, 105K participants). Strongest direct evidence. (TODO: find stable DOI)
\end{enumerate}
\textbf{Synthesis for BehaVoir:} LLMs have the behavioral priors (``effect'') but lack
scientific scaffolding (``channel''). Manning et al.\ is closest to our argument ---
they show LLMs work as scientist AND subject but need structure. BehaVoir provides
that structure via dynamic environment discovery.}

% --- 2026-05-12: Critique-Driven Refinements ---

\idea{Voir = to see; the title is the thesis}{%
``Voir'' means ``to see'' in French. LLMs and current AI operate at the level
of observations --- what can be digitally recorded. Whatever they understand about
us is limited to that, creating a \textbf{jagged view} --- partial, uneven,
observation-bound --- of the space between cognition and behaviorism. BehaVoir
embraces this constraint. The title encodes the thesis: \textbf{BehaVoir} =
behavior, seen. We don't claim to model cognition; we claim that observations
are sufficient for \emph{predictive} behavioral science.}

\idea{Reddit as Atari, not as AlphaGo's Go}{%
Reddit is NOT the contribution --- it is the data substrate. Analogous to Atari
for MuZero: a singular data format hosting a wide array of distinct environments
(subreddits), each with different rules, populations, and reward signals.
MuZero's algorithm operates across 57 Atari games via the same pixel/action
interface; BehaVoir operates across thousands of behavioral environments via
Reddit's uniform structure (posts, comments, votes, timestamps, user histories).
\\[4pt]
Acknowledged limits: population selection (not reflective of world),
content selection (skewed), channel (digital only, ``magic circle'').
But: LLMs \emph{still cannot model even Reddit-level behavior}.
The method requires a singular data format, not specifically Reddit.}

\idea{Anti-circularity principle for LLMs}{%
LLMs trained on Reddit used to reason about Reddit = circular dependency.
Design principle: do \textbf{not} bake knowledge, planning, or policy into
parameters. Use \textbf{written directives and hypotheses} --- explicit,
inspectable, causally verifiable --- for unbiased inference. The LLM is a
reasoning tool and code executor, not a knowledge oracle. Its parametric
knowledge is a \emph{prior} that must be verified against data, never
trusted directly.}

\idea{Epistemic rigor as a core contribution}{%
Humans do science around observational datasets (CMV, etc.). What makes these
structures special? How did humans come up with solving those biases? Theories
developed from these \emph{do} hold in experiments (CMV persuasion findings
replicate in lab). Ideally: maximize real-world experiment validity. Method
discovers dataset $\to$ poses biases/constraints $\to$ fixes end-to-end
$\to$ maximum epistemic rigor $\to$ field experiments hold.
This is not a nice-to-have --- it's a core contribution.}

\idea{Predictive sufficiency, not causal reconstruction}{%
MuZero's deepest lesson: you don't need to reconstruct the true generative
process. Hidden states need not correspond to real states --- they just need
to support accurate predictions. For BehaVoir: we don't need to model cognition.
We need representations that are \emph{predictively sufficient} for behavioral
science. The bar is ``generates hypotheses that hold in field experiments,''
not ``correctly models the human mind.'' Keep this claim soft in the paper ---
let results speak. But it is the philosophical position.}

% --- 2026-05-12: Refined Theoretical Position ---

\idea{BehaVoir learns rules, not representations}{%
The key distinction from MuZero, DreamerV3, and the entire latent world model
lineage: \textbf{BehaVoir's artifacts are explicit hierarchical rules, not latent
states.} We predict \emph{theories} --- inspectable, auditable, causally verifiable
directives that encapsulate an LLM's understanding of a behavioral environment.
\\[4pt]
\textbf{What we share with MuZero/DreamerV3:}
\begin{itemize}
    \item \textbf{Sufficiency}: we don't reconstruct entire behavior to plan and
          understand it. We learn only what's sufficient for scientific decisions.
          Grimm \& Barreto's value equivalence principle (NeurIPS 2020--2022)
          formalizes this: a model need not be accurate about the environment,
          only about what matters for planning. Our rules are the value-equivalent
          model --- sufficient for generating and testing hypotheses.
    \item \textbf{DreamerV3's generality}: one algorithm, fixed hyperparameters,
          150+ tasks. We aim for the same: one method across thousands of Reddit
          environments. DreamerV3 showed larger models $\to$ fewer interactions
          needed. Our scaling aspiration is the same.
    \item \textbf{Planning self-corrects}: MuZero's model is inaccurate but
          compensated by MCTS. Our rules may be imperfect but the empirical
          testing loop (Reddit's time series $\to$ pre-registered experiments)
          corrects them.
\end{itemize}
\textbf{What we differ on --- the critical distinction:}
\begin{itemize}
    \item \textbf{Non-latent artifacts}: MuZero/DreamerV3 learn opaque latent
          dynamics. We learn \emph{explicit rules} --- readable, hierarchical,
          structured. The artifacts themselves are the science. This solves the
          opacity problem that plagues all latent world models.
    \item \textbf{Rules aid LLM understanding}: the hierarchical rules don't
          just compress the environment --- they make the LLM \emph{capable of
          understanding all aspects} available from Reddit for a given question.
          The LLM is a reasoning engine operating on discovered rules, not an
          encoder producing embeddings.
    \item \textbf{Auditability for science}: latent representations cannot be
          peer-reviewed. Rules can. This is not a nice-to-have --- it's what
          makes the output \emph{science} rather than prediction.
\end{itemize}}

\idea{The collapse problem is our noise problem}{%
In the world model literature, the central failure mode is \textbf{representation
collapse} --- the model maps all inputs to the same representation, trivially
satisfying the prediction objective (JEPA's Achilles heel for years; solved only
in 2026 with LeWorldModel's regularization).
\\[4pt]
Our structural analogue: \textbf{environment collapse} --- if the discovery
process lacks structure, all subreddits look equally (un)promising. Raw Reddit
is too noisy to differentiate. Every approach that solved collapse used explicit
regularization:
\begin{itemize}
    \item EfficientZero: self-supervised consistency (does the learned model
          agree with itself over time?)
    \item DreamerV3: symlog normalization, KL balancing, percentile return
          normalization
    \item LeWorldModel: Gaussian regularizer on latent embeddings
\end{itemize}
For BehaVoir, the anti-collapse mechanism is the \textbf{hierarchical rule
structure itself}. Rules that are too generic (``people discuss things'') or
too specific (overfitting to one thread) should be regularized out. The
hierarchy provides the inductive bias: rules at different levels of abstraction
must be internally consistent. This is our analogue of DreamerV3's robustness
techniques.}

\idea{Sufficiency spectrum: our design question}{%
From Grimm \& Barreto and the ``Latent State Design under Sufficiency
Constraints'' line (2025):
\begin{itemize}
    \item \textbf{Observation sufficiency} (reconstruct everything) --- wasteful
    \item \textbf{Prediction sufficiency} (predict future observations) ---
          necessary but not enough for control
    \item \textbf{Control sufficiency} (preserve what affects decisions) ---
          what you actually need
    \item \textbf{Value equivalence} (preserve Bellman-correct value estimates)
          --- tightest useful criterion
\end{itemize}
Critical insight: \textbf{prediction sufficiency does not imply control
sufficiency}. You can perfectly predict what happens next while discarding
distinctions that matter for what to \emph{do}. Applied to BehaVoir: we don't
need rules that predict all Reddit behavior --- we need rules that are
\textbf{sufficient for scientific control}: posing hypotheses, designing
experiments, predicting outcomes. This is a tighter, more useful criterion
than ``understand Reddit.'' And because our rules are explicit (non-latent),
we can \emph{inspect} whether the sufficiency criterion is met.}

% --- 2026-05-12: Undocumented ideas from conversation ---

\idea{The magic circle of digital communication}{%
When people enter Reddit (or any digital platform), they enter a ``magic
circle'' --- a term from game design (Huizinga, 1938) meaning a bounded space
with its own rules, norms, and incentive structures. Subreddits are
\emph{literally} magic circles: explicit rules in sidebars, moderators as
referees, community-specific norms for what counts as a valid contribution.
This is exactly what makes them scientifically useful --- the magic circle
provides the structure that transforms noisy human behavior into something
tractable. BehaVoir's environment discovery is, in essence, \textbf{discovering
magic circles at scale} --- finding bounded behavioral arenas where the rules
are implicit but discoverable.}

\idea{Identity is expressed through communication, not verified externally}{%
On Reddit, ``who'' is revealed through ``what.'' There is no external identity
verification --- users may hide or fake identity. But the key insight is:
\emph{that is the true nature of humans anyway.} People manage impressions in
all social contexts. On Reddit, identity leaks through communication patterns
(MBTI signals from writing style, political leaning from subreddit
participation, expertise from answer quality). This collapses the who/what
distinction: on a text-only platform, identity IS communication.
Methodological implication: we don't need ground-truth identity; behavioral
fingerprints in communication are sufficient.}

\idea{Comments as atomic unit, upvotes as aggregate signal}{%
Reddit provides a dual-signal structure:
\begin{itemize}
    \item \textbf{Comments} = atomic unit of conversation. Rich, granular,
          contextual. Each comment is a behavioral act with author, timestamp,
          content, parent, and position in a thread.
    \item \textbf{Upvotes} = aggregate community signal. We don't have
          individual atomic upvotes (who voted), but the aggregate score
          reflects community judgment. This is a \emph{collective evaluation}
          function --- noisy but massive.
\end{itemize}
This dual structure (granular behavior + aggregate evaluation) is rare in
observational data and is a key reason Reddit works as a behavioral laboratory.}

\idea{The EXPERIGEN $\to$ AlphaZero intermediate step}{%
Between EXPERIGEN (one dataset) and BehaVoir (discover datasets), there is an
intermediate result: the same EXPERIGEN algorithm can be applied to
\textbf{transfer hypotheses across datasets} or \textbf{learn from multiple
sources simultaneously}. This is the AlphaZero step --- same algorithm, multiple
``games,'' no domain-specific modification. This has been observed but is not yet
published. It sits between EXPERIGEN and BehaVoir in the research trajectory
and may become part of this paper or a companion piece.}

\idea{The bottleneck is shared: human AND AI scientists are stalled at datasets}{%
A key framing point: the dataset bottleneck is not just an AI problem. Human
social scientists are \emph{equally} bottlenecked. Experiments are hard to run,
require extensive theorizing, human time, and effort. As data grows, both human
and AI researchers rely more on observational datasets. Progress for
\emph{both} is stalled at the dataset level. This means solving the bottleneck
accelerates the \emph{entire field}, not just the AI pipeline. This should be
a central claim in the introduction.}

\idea{The closed-loop feedback architecture}{%
BehaVoir implies a specific feedback loop, analogous to MuZero's
act--observe--plan cycle:
\begin{enumerate}
    \item \textbf{Discover} an observational environment (subreddit/structure
          with implicit rules)
    \item \textbf{Comprehend} its rules, biases, and constraints
    \item \textbf{Generate} hypotheses about behavior within it
    \item \textbf{Test} via pre-registered experiments (exploit Reddit's
          time-series nature: hypothesize on data before time $T$, validate
          on data after $T$)
    \item \textbf{Refine} both the hypotheses and the environment model
\end{enumerate}
This loop is self-correcting: discovered environments don't need to be perfect,
because the testing step corrects errors (just as MuZero's MCTS compensates
for dynamics model inaccuracies). The loop IS the method.}

\idea{Three levels of contribution (from critique analysis)}{%
Cleanly separate what is novel:
\begin{enumerate}
    \item \textbf{Environment discovery}: Automate identification of clean
          observational environments from raw data. This is the ``What +
          Context'' problem humans have done manually for decades. [NEW]
    \item \textbf{Environment comprehension}: Learn the implicit rules, norms,
          biases, and reward structures of discovered environments. [NEW]
    \item \textbf{Hypothesis generation and testing}: Apply EXPERIGEN-style
          hypothesis search within discovered environments. [INHERITED from
          EXPERIGEN, but now at scale across many environments]
\end{enumerate}
Levels 1 and 2 are BehaVoir's novel contributions. Level 3 is inherited
but demonstrated at unprecedented scale.}

\idea{The narrative arc of the paper}{%
Introduction should follow the DeepMind trajectory as a structural device:
\begin{enumerate}
    \item Social science is slow. Iterative cycles of theorize--experiment--validate.
    \item EXPERIGEN proved AI can do end-to-end social science on curated data
          (the AlphaGo moment).
    \item But curated data is the bottleneck. Only a handful of datasets exist,
          each taking years of human effort to identify and validate.
    \item This bottleneck stalls \emph{both} human and AI scientists.
    \item BehaVoir removes it: dynamically discover observational environments
          from raw data, comprehend their structure, and run science at scale.
    \item Reddit is the first platform (Atari) --- singular format, many environments.
    \item The method generalizes beyond Reddit to any uniformly structured
          behavioral data source.
\end{enumerate}}

\idea{MuZero analogy must be stated precisely}{%
Readers will misread the MuZero analogy as ``discovering new environments.''
But MuZero's contribution is learning a dynamics model \emph{within} a single
environment without being given the rules. The correct mapping:
\begin{itemize}
    \item EXPERIGEN needs a pre-curated dataset = AlphaZero needs game rules
    \item BehaVoir learns the structure from raw observations = MuZero learns
          the dynamics from raw pixels
\end{itemize}
MuZero doesn't discover which game to play --- it removes the dependency on
being \emph{given} the rules. Similarly, BehaVoir removes the dependency on
being \emph{given} curated datasets. State this explicitly in the paper or
reviewers will attack it.}

\idea{``We are here'' on the MuZero diagram}{%
User referenced the MuZero blog post diagram
(\texttt{deepmind.google/blog/muzero-mastering-go-chess-shogi-and-atari-without-rules/}).
The positioning is: we are at the point where the environment model is being
\emph{learned}, not just used. EXPERIGEN assumed a clean environment (curated
dataset); BehaVoir learns what the environment IS from raw observations.
Consider including a version of this diagram (adapted for our trajectory)
in the introduction or related work.}

% --- 2026-05-12: From paper reading + critique session ---

\idea{EXPERIGEN's limitations define BehaVoir's scope}{%
From reading the EXPERIGEN paper, these are the specific walls it hits:
\begin{enumerate}
    \item \textbf{One dataset at a time}: no transfer across domains. The
          hypothesis bank is local to a single dataset.
    \item \textbf{No unified model}: produces independent hypotheses, never
          synthesized into a coherent theory.
    \item \textbf{No causal reasoning}: all results are correlational/observational.
          Explicitly avoids causal claims.
    \item \textbf{LLM inference plateaus at ${\sim}20$ outer iterations}:
          context window saturation degrades hypothesis quality. AutoML variant
          continues improving to 40 iterations --- suggesting a hybrid handoff.
    \item \textbf{Feature extractor noise sensitivity}: even 10--20\% label error
          degrades both accuracy and hypothesis count.
    \item \textbf{No multi-variable hypotheses}: generates individual hypotheses,
          not conditional or interactive ones (though iterative refinement helps).
\end{enumerate}
BehaVoir inherits EXPERIGEN's hypothesis engine (level 3 contribution) but must
overcome these limits through environment discovery (level 1) and comprehension
(level 2). The AutoML handoff insight is particularly relevant: LLMs for
interpretable discovery, AutoML for scaling.}

\idea{Superhuman through information asymmetry}{%
The argument for superhuman performance is not just compute --- it's
\textbf{information asymmetry}. No individual human expert possesses knowledge
equivalent to the entirety of Reddit. An expert in persuasion knows CMV deeply
but not r/relationship\_advice. A mental health researcher knows r/SuicideWatch
but not r/AITA's moral reasoning patterns. BehaVoir, operating across all
environments simultaneously, has access to cross-domain patterns no single human
can hold. Superhuman here means: discovering behavioral regularities that span
environments in ways no expert's training prepares them to see. We will
benchmark against experts at the end and let the comparison define the claim.}

\idea{Observational findings DO replicate in controlled experiments}{%
A crucial empirical observation the user emphasized: theories developed from
observational datasets (CMV persuasion findings, mental health indicators,
moral reasoning patterns) \emph{do come out to be true in controlled
experiments}. This is not obvious --- the entire field of causal inference
exists because observational data is supposed to be unreliable. But the
structures these datasets provide (clean formulation, natural labels, community
norms) apparently create conditions that approximate experimental control.
This is BehaVoir's epistemic foundation: if we can discover environments with
similar structural properties, their observational findings should also
replicate. The goal is to maximize this: method discovers dataset $\to$ poses
biases/constraints $\to$ \emph{fixes} them end-to-end $\to$ field experiments
hold.}

\idea{Learning WHAT to learn, not just predicting}{%
The user's phrasing was precise: ``the key decisions that improve its
understanding.'' This is more active than plain predictive sufficiency.
The system doesn't just predict --- it learns \emph{what to learn next}.
Which environment to explore? Which bias to address? Which hypothesis to
test first? This is the planning component (MuZero's MCTS analogue):
making decisions about where to direct scientific attention for maximum
epistemic gain. The system learns a policy over scientific decisions,
not just a model of behavior.}

\idea{Dreamer's sample efficiency lesson}{%
From reading the Dreamer paper: it achieves ${\sim}20{\times}$ fewer
environment steps than model-free baselines by learning in imagination rather
than real experience. For BehaVoir: ``real experiments'' (actually testing
hypotheses on live data) are expensive. If we can learn a sufficiently good
environment model, we can generate and pre-filter hypotheses in ``imagination''
(simulated environments) before spending compute on real validation. This
is the sample efficiency argument: fewer pre-registered experiments needed
per unit of scientific progress.}

\idea{Soft-claim strategy: let results prove the thesis}{%
Explicit writing strategy: keep the ``communication is sufficient'' claim soft
in the paper. Don't make it a hard philosophical stance. Instead, frame it as:
``we explore what can be achieved with observations alone'' and let the
experimental results demonstrate the reach and limits. This avoids getting
bogged down in cognition-vs-behaviorism debates and keeps the paper focused
on the empirical contribution. The title (BehaVoir = behavior, seen) makes
the position clear without requiring a philosophical defense.}

%  RAW SNIPPETS / HALF-THOUGHTS
\section{half thoughts}

% - "slowly we..." — the trajectory from EXPERIGEN to BehaVoir is gradual:
%   one-dataset discovery → multi-dataset transfer → environment discovery.
%   Each step removes a dependency. The paper should feel like a natural
%   next step, not a discontinuous leap.
%
% - The field naming question is settled (CSS) but the paper title uses
%   "Human Behavior" not "Social Science" — this is intentional. BehaVoir
%   is about behavior specifically, not all of social science. The narrowing
%   is strategic: behavior can be observed; institutions, policies, markets
%   are harder to observe purely from communication.
%
% - Reddit's magic circle has a temporal dimension too: subreddit rules
%   change over time, moderator policies shift, community norms evolve.
%   A discovered environment is a snapshot. The time-series nature lets
%   us track how environments themselves change — meta-discovery.
%
% - The anti-circularity principle (written directives, not baked parameters)
%   connects to a broader philosophy: the system should be auditable. Every
%   hypothesis, every environmental judgment, every bias assessment should
%   be traceable back to explicit reasoning, not opaque model internals.
%   This is not just about avoiding circularity — it's about epistemic
%   transparency as a first-class concern.
%
% - Reddit covers most topics that get discussed, but with skew. This is
%   different from "content selection bias" — it's more like a wide-angle
%   lens with distortion. The coverage is vast; the faithfulness varies.
%
% - "Between cognition and behaviorism" — BehaVoir is not squarely in
%   either camp. It's observation-only (behaviorist in method) but aims to
%   understand behavior (not just predict it). The jagged view is the space
%   in between: what you can see of the mind through actions alone. This is
%   a novel epistemological position worth articulating carefully.
%
% - EXPERIGEN's A/B test result (+344% uplift, p < 1e-6) is powerful
%   ammunition for the introduction. It proves the hypothesis engine works
%   in the real world. BehaVoir's claim is: now apply that engine at scale,
%   across dynamically discovered environments.
%
% - The AutoML handoff from EXPERIGEN (LLMs plateau at ~20 iterations,
%   AutoML continues to 40) suggests BehaVoir may need a similar hybrid:
%   LLMs for interpretable environment discovery, automated methods for
%   scaling hypothesis testing within discovered environments.
% BehaVoir — Todo Tracker
% Use \pending{} for open items, \done{} for completed
% Use \experiment{name}{status} for experiment tracking

\section*{Todo}

\subsection*{Framing \& Scope}
\done{Clarify field name --- use \textbf{Computational Social Science} as primary; ``AI for Social Science'' as secondary for ML audience. See scratch.tex for full rationale.}
\done{Compile list of domains where Reddit is used in AI/Social Science --- 12 Reddit + 3 non-Reddit datasets cataloged in scratch.tex with papers, sizes, and why each is ``clean''}

\subsection*{Writing}
\pending{Write Introduction — frame the AlphaGo $\to$ AlphaZero $\to$ MuZero analogy}
\pending{Related Work — observational studies, automated science, dataset curation}
\pending{Methodology section outline}

\subsection*{Literature \& References}
\done{Survey behavior simulation with LLMs --- 7 key papers added to references.bib and scratch.tex (Park, Horton, Manning, Aher, Argyle, Santurkar, Hewitt). Hewitt et al.\ still needs stable DOI.}
\done{Cite Lasswell's communication model (1948) and position our mapping --- added \texttt{lasswell1948structure}, \texttt{buntain2014identifying} to references.bib}

\subsection*{Methodology \& Theory}
\pending{Define pre-registration protocol: formulate hypotheses on data up to time $T$, freeze, test on data after $T$. Specify precisely --- cannot just assert ``Reddit is a time series.''}
\pending{Solve the epistemic rigor bottleneck: how did humans solve biases in observational datasets (CMV, etc.)? What makes these structures special? How to maximize real-world experiment validity --- method discovers dataset, poses biases/constraints, fixes end-to-end, so field experiments hold.}
\pending{Operationalize ``superhuman performance'' --- define metrics for comparison with human experts.}

\section{Vidit's Rough}

\subsection{Subreddit Selection Rationale}

\begin{table}[h]
\centering
\small
\begin{tabular}{|l|l|}
\hline
\textbf{Subreddit} & \textbf{Why Selected} \\
\hline
\textbf{nostalgia} & Dedicated memorability/nostalgia community \\
\textbf{tipofmytongue} & People searching for half-remembered media \\
\textbf{tipofmyjoystick} & Same pattern for games \\
\textbf{Glitch\_in\_the\_Matrix} & Vivid anomalous experiences \\
\textbf{MandelaEffect} & Collective false memories \\
\textbf{OldSchoolCool} & Visual nostalgia, memorable past images \\
\textbf{bestof} & Community-curated memorable content \\
\textbf{DoesAnybodyElse} & Shared memorable experiences/feelings \\
\textbf{AskOldPeople} & Older adults reflecting on memorable events \\
\textbf{CasualConversation} & Personal anecdotes about memorable moments \\
\textbf{TrueOffMyChest} & Confessions about memorable events \\
\textbf{offmychest} & Personal stories memorable enough to share \\
\textbf{tifu} & Memorable embarrassing/disastrous stories \\
\textbf{confession} & Confessions of memorable acts/events \\
\textbf{AskReddit} & Many ``what's the most memorable X'' threads (title-filtered) \\
\textbf{movies} & ``Most memorable scene/film'' discussions (title-filtered) \\
\textbf{books} & ``Books that stayed with you'' discussions (title-filtered) \\
\textbf{gaming} & ``Most memorable gaming moment'' (title-filtered) \\
\textbf{Games} & Discussion-focused gaming (title-filtered) \\
\textbf{Music} & ``Songs you'll never forget'' (title-filtered) \\
\textbf{travel} & Memorable travel experiences (title-filtered) \\
\textbf{todayilearned} & Memorable facts (title-filtered) \\
\textbf{Showerthoughts} & Observations about memory (title-filtered) \\
\textbf{relationships} & Memorable relationship moments (title-filtered) \\
\textbf{relationship\_advice} & Memorable relationship events (title-filtered) \\
\hline
\end{tabular}
\end{table}

\subsubsection*{Rejected:}
\begin{itemize}
    \item r/2000sNostalgia (26 shards), r/gamenostalgia (14 shards): minimal corpus presence
    \item r/AustralianNostalgia, r/IndiaNostalgia: regional, not generalizable
    \item r/dreams, r/WritingPrompts: fiction/niche, memorability incidental
    \item r/pics, r/mildlyinteresting: visual content, comment text doesn't discuss memorability
\end{itemize}

\subsubsection{Corpus Limitations}
\begin{enumerate}
    \item \textbf{Temporal skew}: March--April 2020 = early COVID-19 lockdowns. Memorability discussions biased toward pandemic experiences.
    \item \textbf{Missing submissions}: 25.3\% of threads have no submission in the corpus (likely deleted/removed posts). These threads are kept if comments show memorability signals.
    \item \textbf{AskReddit sampling}: Randomly sampled to 30\% cap; some memorability threads may have been excluded.
    \item \textbf{Weak supervision precision}: Lexical LFs miss implicit memorability expressed without standard phrases.
\end{enumerate}

\subsubsection{Alternatives Considered and Rejected}
\begin{enumerate}
    \item \textbf{Embedding-based semantic similarity}: Better for nuanced memorability but risks timeout at scale.
    \item \textbf{Blanket Tier 1 positive labelling (iteration 1 approach)}: Produced \textasciitilde40\% false positive rate; replaced with uniform weak supervision.
    \item \textbf{Score-based stratified sampling}: Research question is about memorability regardless of engagement.
    \item \textbf{LLM-based labelling}: Gold standard quality but infeasible at scale.
    \item \textbf{Tier 1 only collection}: High precision but misses cross-discourse signals in r/gaming, r/movies, etc.
\end{enumerate}


\subsection*{Experiments}
% \experiment{name}{status}


\bibliographystyle{plain}
\bibliography{references}

\end{document}
